\documentclass[11pt,a4paper]{article}

\usepackage[utf8]{inputenc}
\usepackage{amsmath,amssymb,amsthm}
\usepackage{geometry}
\usepackage{hyperref}
\usepackage{booktabs}
\usepackage{xcolor}
\usepackage{enumitem}
\usepackage{url}

\geometry{margin=2.5cm}
\hypersetup{colorlinks=true,linkcolor=blue,citecolor=blue,urlcolor=blue}

\title{Design Document: RL-Driven Physical Activity Interventions\\
\large A Configurable Simulation Framework}
\author{William Dennis\\
Bristol x NUS Research Project\\
\texttt{ky23812@bristol.ac.uk}}
\date{June 2026}

\begin{document}
\maketitle
\tableofcontents

\newpage

\section{Introduction}
\label{sec:intro}

Reinforcement learning offers a principled approach to personalising
just-in-time adaptive interventions (JITAIs) for physical activity, but
training RL agents requires either expensive real-world trials or a realistic
simulation environment. Existing simulators are bespoke and tied to specific
datasets or papers, making cross-paper comparison difficult and limiting
reproducibility. We present \texttt{rl-health-interventions}, a configurable
simulation framework where the dataset schema, MDP dynamics, user behaviour
models, and agent configurations are specified entirely through YAML config
files---no source code changes required. The framework enables systematic
comparison of RL agents on PA intervention tasks using synthetic data in
Phase 1, with real-data integration planned for Phase 2.

The initial implementation delivers a rule-based user simulator with four
behavioural archetypes, a multi-timescale MDP environment (daily steps +
delayed 3-week clinical measures), and Thompson Sampling as the baseline
agent \cite{karine2024stepcountjitai,klasnja2019heartsteps}. Deep RL agents
(DQN, PPO) will be added if time permits. All
components share a common ABC + registry interface for extensibility.

\section{MDP Formalisation}
\label{sec:mdp}

This section specifies the Markov Decision Process (MDP) for the first
iteration of the simulation framework. The MDP is defined as the tuple
$(\mathcal{S}, \mathcal{A}, P, R, \gamma)$.

\subsection{State Space $\mathcal{S}$}

The state represents the current context of a user at decision epoch $t$.
States are constructed from wearable-device signals, user profile
information, and intervention history:

\[
s_t = \{ \mathbf{x}_t^{\text{wearable}}, \mathbf{x}_t^{\text{context}},
\mathbf{x}_t^{\text{history}}, \mathbf{x}^{\text{profile}} \}
\]

\begin{table}[h]
\centering
\caption{State variables}
\begin{tabular}{lll}
\toprule
Variable & Type & Description \\
\midrule
$\text{steps}_t$ & $\mathbb{R}_{\geq 0}$ & Step count, past 24h \\
$\text{hr}_t$ & $\mathbb{R}_{\geq 0}$ & Heart rate (or resting) \\
$\text{sleep\_hours}_t$ & $\mathbb{R}_{\geq 0}$ & Previous night sleep \\
$\text{sedentary\_min}_t$ & $\mathbb{R}_{\geq 0}$ & Sedentary minutes today \\
$\text{time\_of\_day}_t$ & $\{0,\dots,23\}$ & Current hour \\
$\text{day\_of\_week}_t$ & $\{0,\dots,6\}$ & Day of week \\
$\text{goal\_progress}_t$ & $[0,1]$ & Fraction of daily goal met \\
$\text{burden}_t$ & $\mathbb{Z}_{\geq 0}$ & Interventions sent today \\
$a_{t-1}$ & $\mathcal{A}$ & Previous action \\
$\text{response}_{t-1}$ & $\{0,1\}$ & Previous response \\
$\text{body\_measure}_k$ & $\mathbb{R}$ & Clinical measure (q3wks) \\
$\text{age}$ & $\mathbb{Z}_{>0}$ & Static profile \\
$\text{gender}$ & $\{0,1\}$ & Static profile \\
$\text{baseline\_activity}$ & $\{0,1,2\}$ & Low / medium / high \\
\bottomrule
\end{tabular}
\end{table}

Variables are normalised to $[0,1]$ before being passed to the agent. The
body measure is only observed every 3 weeks (sparse reward signal).

\subsection{Action Space $\mathcal{A}$}

\[
\mathcal{A} = \{ a_0, a_1, a_2, a_3, a_4, a_5 \}
\]

\begin{table}[h]
\centering
\caption{Action space}
\begin{tabular}{cll}
\toprule
Action & Label & Description \\
\midrule
$a_0$ & No message & Do not intervene \\
$a_1$ & Motivational prompt & Generic encouragement \\
$a_2$ & Walking suggestion & Recommend a short walk \\
$a_3$ & Goal reminder & Remind of step goal \\
$a_4$ & Recovery suggestion & Stretch or rest \\
$a_5$ & Progress feedback & Report goal progress \\
\bottomrule
\end{tabular}
\end{table}

Exactly one action per decision epoch. The ``no message'' action is critical
for managing notification burden.

\subsection{Transition Function $P$}

The transition function models the probability over next states:

\[
P(s_{t+1} \mid s_t, a_t): \mathcal{S} \times \mathcal{A} \times \mathcal{S}
\rightarrow [0,1]
\]

In Phase 1 (rule-based), transitions follow a behavioural response model:

\[
\Delta \text{steps}_{t+1} = f(s_t, a_t, \theta_{\text{user}})
\]

where $\theta_{\text{user}}$ parameterises one of four behavioural archetypes:

\begin{itemize}[noitemsep]
    \item \textbf{Goal-driven:} Responds to reminders and feedback.
    Proportional to goal progress.
    \item \textbf{Social responder:} Responds to motivational prompts.
    Models social desirability bias.
    \item \textbf{Resistant:} Low response. Fatigue accumulates quickly.
    \item \textbf{Stable maintainer:} Already active. Low marginal gain.
\end{itemize}

Burden evolves as:

\[
\text{burden}_{t+1} =
\begin{cases}
\text{burden}_t + 1 & \text{if } a_t \neq a_0 \\
0 & \text{if } a_t = a_0
\end{cases}
\]

When $\text{burden} > B_{\max}$, response probability decays linearly.

\subsection{Reward Function $R$}
\label{sec:reward}

\subsubsection{Immediate Reward}

\[
R_t^{\text{immediate}} = \alpha \cdot \Delta\text{steps}_t
- \beta \cdot \mathbf{1}_{\{a_t \neq a_0\}}
+ \lambda \cdot \text{goal\_progress}_t
\]

\subsubsection{Delayed Reward (Body Measures)}

Clinical body measures arrive every 3 weeks ($t = 3k$):

\[
R_{3k}^{\text{delayed}} = \eta \cdot \text{BM\_improvement}_{3k}
\]

Total reward:

\[
R_t =
\begin{cases}
R_t^{\text{immediate}} + R_t^{\text{delayed}} & t \equiv 0 \pmod{3} \\
R_t^{\text{immediate}} & \text{otherwise}
\end{cases}
\]

The compound reward (immediate + delayed) is a central research question.
Weights $(\alpha, \beta, \lambda, \eta)$ are configurable experiment parameters.

\subsection{Discount Factor $\gamma$}

\[
\gamma \in [0.9, 0.95]
\]

A high discount factor encourages long-term behaviour change over short-term
step increases. Appropriate for health habit formation.

\subsection{Decision Epoch}

\textbf{Daily} decisions. This aligns with typical intervention cadence,
the 3-week body measure delay (3 epochs), and practical deployment
constraints.

\section{Decision Log}
\label{sec:decisions}

\subsection{Confirmed Decisions}

\begin{table}[h]
\centering
\begin{tabular}{cll}
\toprule
\# & Decision & Status \\
\midrule
1 & Config format: YAML & Confirmed \\
2 & Interface: CLI + config files; web UI = stretch & Confirmed \\
3 & Python 3.11 + \texttt{uv} for package management & Confirmed \\
4 & Synthetic data for Phase 1 (real data requires 4-8 week applications) & Confirmed \\
5 & Thompson Sampling as initial agent & Confirmed \\
6 & Daily decision epochs & Confirmed \\
7 & Multi-timescale reward: immediate (daily steps) + delayed (3-week body measures) & Confirmed \\
8 & Four behavioural archetypes: goal-driven, social responder, resistant, stable maintainer & Confirmed \\
9 & Pluggable component architecture (ABC + registry pattern) & Confirmed \\
10 & Minimal dependencies: no Gymnasium, no SB3 & Confirmed \\
\bottomrule
\end{tabular}
\end{table}

\subsection{Open Decisions}

The following require supervisor input:

\begin{description}[noitemsep]
    \item[Distribution model] \hfill \\
    Installable package (\texttt{uv build} produces a wheel) vs repo-based
    (clone and \texttt{uv sync})?

    \item[MDP baseline approval] \hfill \\
    Does the MDP formalisation (\S2) look like a reasonable starting point?

    \item[Phase 2 dataset targets] \hfill \\
    HeartSteps first (fastest access), then All of Us, then UK Biobank?

    \item[Results table format] \hfill \\
    What format and content do you want in the experiment output (CSV,
    terminal table, JSON, summary plots)?

    \item[Simulator validation criteria] \hfill \\
    Acceptable thresholds for response-rate RMSE, state-distribution KL
    divergence, and engagement-curve alignment before running experiments?
\end{description}

\section{Future Work}

Beyond Phase 1, the framework targets LLM-based user simulation using
prompt-driven persona modelling, offline RL agents (CQL, IQL) for batch
policy learning, non-stationary environment adaptation, and real-data
validation against All of Us, UK Biobank, and HeartSteps datasets.

\bibliographystyle{unsrt}
\bibliography{references}

\end{document}
